\documentclass[11pt]{letter}
\usepackage[margin=1in]{geometry}
\usepackage{hyperref}

\signature{Ian Todd\\Sydney Medical School\\University of Sydney\\itod2305@uni.sydney.edu.au}
\address{Ian Todd\\Sydney Medical School\\University of Sydney\\Sydney, NSW 2006\\Australia}

\begin{document}

\begin{letter}{The Editors\\Physical Review E\\American Physical Society}

\opening{Dear Editors,}

I am pleased to submit our manuscript, ``The Dimensional Work of Surfaces: A Thermodynamic Origin of Minimal Geometry and Surface Tension,'' for consideration as a Research Article in \textit{Physical Review E}.

\textbf{Summary.} Surface tension and minimal surfaces are foundational concepts in soft matter physics, geometry, and biological interfaces. Classical theory describes these phenomena through postulated energy functionals---$\gamma A$, the Young--Laplace law, and Helfrich curvature elasticity---but the thermodynamic origin of these geometric energies has remained opaque.

In this work, we show that classical surface energetics emerge naturally from the \emph{Dimensional Landauer Bound}, a recently introduced generalization of Landauer's principle to high-dimensional stochastic systems constrained to lower-dimensional manifolds. When a noisy physical system is confined to a two-dimensional surface embedded in three dimensions, the minimum work required to impose that constraint is given by a KL divergence that, in the continuum limit, reduces to a functional of the induced metric:
\[
W_{\mathrm{dim}} \propto k_B T \int \log\sqrt{\det G(u,v)}\, p(u,v)\,du\,dv.
\]
Expanding this functional yields the area functional, the Young--Laplace pressure, and Helfrich curvature elasticity as successive orders of dimensional work.

\textbf{Key Results.}
\begin{enumerate}
\item Surface tension $\gamma A$ emerges as the leading-order dimensional work for constraining fluctuations to a curved manifold.
\item The Young--Laplace equation $\Delta P = 2\gamma H$ arises from variational minimization under volume constraint.
\item Mean curvature flow is the gradient descent of dimensional work, with minimal surfaces as fixed points.
\item Numerical simulations confirm that relaxation toward minimal surfaces corresponds to dimensional work minimization.
\end{enumerate}

\textbf{Why Physical Review E?} This manuscript bridges stochastic thermodynamics, soft matter physics, and differential geometry---all core areas of PRE's scope. We provide a unified thermodynamic derivation of surface energetics from information-theoretic principles, connecting classical capillarity to modern work on fluctuation theorems and nonequilibrium thermodynamics. The result that surface tension and curvature elasticity both arise from a single underlying principle---dimensional contraction---represents a conceptual advance that should interest the PRE readership across statistical physics, soft matter, and biological physics.

We confirm that this manuscript has not been published and is not under consideration elsewhere. We have no conflicts of interest to declare.

\closing{Sincerely,}

\end{letter}
\end{document}
